% =============================================================================
% Section 1 — Introduction
% =============================================================================
\section{Introduction}
\label{sec:intro}

Modern distributed applications routinely require storage systems that remain
available during network partitions, tolerate node failures without service
interruption, and scale horizontally without bottlenecks imposed by a central
coordinator. The CAP theorem~\cite{brewer2000cap,gilbert2002cap} establishes
that no distributed data store can simultaneously guarantee strong consistency,
availability, and partition tolerance. Systems that choose availability and
partition tolerance---the AP corner of the CAP triangle---must accept that nodes
may temporarily diverge and must define a reconciliation strategy for concurrent
writes.

Consensus-based systems such as Raft~\cite{ongaro2014raft} and
Paxos~\cite{lamport2001paxos} achieve strong consistency by electing a leader
that serialises all writes, at the cost of availability when a quorum cannot be
reached, and of increased operational complexity. Leaderless, eventually consistent
systems such as Amazon Dynamo~\cite{decandia2007dynamo} and Apache Cassandra
sacrifice linearisability in favour of write availability on every node and
partition tolerance, using vector clocks or timestamps to reconcile diverged state.

\pebble{} occupies this leaderless, eventually consistent space. It is a
distributed key-value store built on a small set of composable primitives:

\begin{itemize}[leftmargin=1.5em]
  \item \textbf{Immutable, signed mutations.} Every write or delete is captured
        as a cryptographically signed, globally unique record that is never
        modified after creation.
  \item \textbf{Deterministic Last-Write-Wins (LWW) conflict resolution.} A
        total order over mutations resolves concurrent writes without coordination,
        ensuring that all nodes independently compute the same winning value.
  \item \textbf{Pull-based gossip replication.} Nodes periodically exchange
        version summaries and transfer only the mutations each peer is missing,
        providing bandwidth-efficient eventual convergence.
  \item \textbf{Ed25519 authentication.} Both inter-node connections and external
        client sessions are authenticated with digital signatures, preventing
        rogue nodes or forged mutations from corrupting cluster state.
\end{itemize}

\subsection{Design Goals}

\pebble{} is designed to satisfy the following goals:

\begin{description}[leftmargin=1.5em, labelwidth=3.5cm]
  \item[Leaderless.] Any node accepts reads and writes. There is no leader
        election and no single point of failure.
  \item[Eventually consistent.] All healthy nodes converge to identical state
        after the last write, given sufficient gossip rounds.
  \item[Authenticated.] Mutations carry Ed25519 signatures; nodes and clients
        are verified before any state exchange occurs.
  \item[Simple.] The core mechanisms (LWW, gossip, version vectors) are well
        understood and compose cleanly. The entire system is implemented in
        fewer than 3,000 lines of TypeScript.
  \item[Durable.] Local mutations are written to an append-only log and synced
        to disk before acknowledgment. Nodes recover from crash by replaying the
        log from the last snapshot.
\end{description}

\subsection{Non-Goals (v0.1)}

The following properties are explicitly \emph{not} provided in the current
version and are discussed as future work in \cref{sec:future}:

\begin{itemize}[leftmargin=1.5em]
  \item Linearisability or sequential consistency.
  \item Multi-key atomic transactions.
  \item Dynamic cluster membership (node join/leave without restart).
  \item Transport-layer encryption (TLS).
  \item Bounded divergence guarantees under arbitrary clock skew.
\end{itemize}

\subsection{Contributions}

This whitepaper makes the following technical contributions:

\begin{enumerate}[leftmargin=1.5em]
  \item A formal definition of the mutation algebra and the LWW total order, with
        a proof of convergence (\cref{sec:conflict}).
  \item A precise specification of the version-vector-based incremental
        synchronisation scheme (\cref{sec:vv}).
  \item A complete description of the 4-step symmetric gossip protocol, with
        deadlock-freedom and idempotency arguments (\cref{sec:gossip}).
  \item A complexity analysis of per-round bandwidth, state growth, and
        cluster-wide convergence time as a function of node count $N$
        (\cref{sec:complexity}).
  \item A catalogue of eight formally stated system invariants (\cref{sec:invariants}).
\end{enumerate}

\subsection{Paper Organisation}

\Cref{sec:arch} describes the overall system architecture.
\Cref{sec:datamodel} defines the data model.
\Cref{sec:crypto} covers the cryptographic layer.
\Cref{sec:conflict} presents the LWW conflict resolution mechanism.
\Cref{sec:vv} formalises version vectors.
\Cref{sec:gossip} specifies the gossip protocol.
\Cref{sec:nodeauth,sec:clientauth} describe node and client authentication.
\Cref{sec:wire} details TCP framing.
\Cref{sec:persistence} covers persistence and crash recovery.
\Cref{sec:complexity} provides complexity analysis.
\Cref{sec:invariants} enumerates system invariants.
\Cref{sec:limitations,sec:future} discuss limitations and future work.
\Cref{sec:conclusion} concludes.
